\chapter[Descenso de los morfismos étales]{Descenso de los morfismos étales.\\ Aplicación al grupo fundamental}
\label{IX}
\marginpar{228}

\section{Recordatorio sobre los morfismos étales}
\label{IX.1}

Vamos a repasar aquí las propiedades de los morfismos étales (desarrolladas
en~\Ref{I}) que vamos a utilizar, aprovechando la ocasión para eliminar de
la teoría las hipótesis noetherianas superfluas. El lector observará que,
aun cuando solo interesen los esquemas noetherianos, la técnica del descenso
lleva a introducir esquemas no noetherianos (tales como
$\Spec(\widehat{A}\otimes_{A}\widehat{A})$, donde~$A$ es un anillo local
noetheriano); para poder aplicar el lenguaje de las categorías fibradas,
importa definir nociones como la de morfismo étale,
\ifthenelse{\boolean{orig}}{etc...}{etc.}, sin introducir en ellas ninguna
restricción noetheriana.
El lector que no desee comprobar ni admitir que los enunciados siguientes
son verdaderos sin hipótesis noetherianas podrá limitarse a admitirlos bajo
las hipótesis noetherianas de~\Ref{I}, con tal de introducir esas mismas
hipótesis noetherianas en los enunciados de los números siguientes y de
utilizar la definición~\Ref{IX.1.1} que se da a continuación para los
esquemas no noetherianos que aparecen en los razonamientos.

\begin{definition}
\label{IX.1.1}
Sean~$f\colon X\to S$ un morfismo de preesquemas y~$x$ un punto de~$X$.
Se dice que~$f$ es \emph{étale en}~$x$,
\index{étale (morfismo, álgebra)|hyperpage}%
o que~$X$ es étale sobre~$S$ en~$x$, si existen un entorno abierto
afín~$U$ de~$s=f(x)$, un entorno abierto afín~$V$ de~$x$ situado sobre~$U$,
un esquema afín \emph{noetheriano}~$U_{0}$, un~$U_{0}$-esquema
afín~$V_{0}$, \emph{étale} en el sentido de~\eqref{I}, un
morfismo~$U\to U_{0}$ y un
$U$-isomorfismo
$$
V\isomto V_{0}\times_{U_{0}}U\quoi.
$$
\end{definition}

Se observará
\marginpar{229}
que, cuando~$S$ es localmente noetheriano, esta terminología coincide con
la de~\loccit Se dirá asimismo que~$f$ \emph{es étale}, o que~$X$
\emph{es étale sobre}~$S$, si~$f$ es étale en todo punto~$x$ de~$X$.
Con estas definiciones, las proposiciones siguientes se reducen sin
dificultad al caso noetheriano, donde se demuestran en~\Ref{I.4},
\Ref{I.5} y~\Ref{I.7}. Para más detalles, el lector podrá consultar
EGA~IV\footnote{Más precisamente, EGA~IV 17,~18.}.

\begin{remarques}
\label{IX.1.2}
Si~$f$ es étale en~$x$, entonces~$f$ es «\emph{de presentación finita
en}~$x$» (\Ref{VIII.3.5}), el anillo local de~$x$ en la fibra
$f^{-1}(s)$ es una \emph{extensión finita separable} de~$\kres(s)$ y,
finalmente, $f$ es \emph{plano en}~$x$. Puede demostrarse que la recíproca
es cierta; por tanto, la definición~\Ref{IX.1.1} es la misma que cuando~$S$
es localmente noetheriano, salvo que debe reemplazarse la condición «de
tipo finito en~$x$» por «de presentación finita en~$x$». Como la
demostración de este resultado es delicada, no hemos querido adoptar aquí
esta definición de la noción de morfismo étale, que no se presta
directamente a la demostración de las propiedades que siguen.
\end{remarques}

Observemos primero que se tiene trivialmente:
\begin{proposition}
\label{IX.1.3}
Si~$f\colon X\to S$ es étale, entonces todo morfismo~$f'\colon X'\to S'$
que se deduce de él mediante un cambio de base~$S'\to S$ es también étale.
\end{proposition}

Puede decirse, por tanto, que los morfismos étales forman una
\emph{subcategoría fibrada} de la categoría de las flechas en~$\Sch$
\ifthenelse{\boolean{orig}}
{(\cf \Ref{VI.11}~a)}
{(\cf \Ref{VI.11}~a))}. El objeto de la presente exposición es estudiar
las propiedades de exactitud de esta categoría fibrada sobre~$\Sch$.
\begin{proposition}
\label{IX.1.4}
Sea~$f\colon X\to S$ un morfismo de preesquemas. Para que sea una inmersión
abierta es necesario y suficiente que sea étale y radicial.
\end{proposition}

\Cf \Ref{I.5.1}. Se deduce que, si~$X$ es étale sobre~$S$, toda sección
de~$X$ sobre~$S$ es una inmersión abierta; utilizando de nuevo~\Ref{IX.1.4},
se obtiene:
\begin{corollaire}
\label{IX.1.5}
Sea~$X$ un $S$-preesquema étale. Entonces existe una correspondencia
biunívoca entre el conjunto de las secciones de~$X$ sobre~$S$ y el conjunto
de las partes
\marginpar{230}
abiertas~$\Gamma$ de~$X$ tales que el morfismo~$\Gamma\to S$ inducido por
el morfismo estructural sea \emph{radicial} y \emph{sobreyectivo}.
\end{corollaire}

Si, además, $X$ está separado sobre~$S$, $\Gamma$ será una parte de~$X$ a la
vez abierta y cerrada, pero esto poco importa. --- Efectuando un cambio de
base evidente, puede ponerse~\Ref{IX.1.5} bajo la forma aparentemente más
general siguiente:
\begin{corollaire}
\label{IX.1.6}
Sean~$X$ e~$Y$ dos $S$-preesquemas, siendo~$Y$ étale sobre~$S$. Entonces la
aplicación~$f\mto\Gamma_f$ que asocia a todo $S$-morfismo~$f$ de~$X$ en~$Y$
la parte de~$X\times_S Y$ subyacente al grafo de~$f$ es una biyección de
$\Hom_S(X,Y)$ sobre
\ifthenelse{\boolean{orig}}
{el conjunto de las partes abiertas de~$\Gamma$ de~$X\times_S Y$}
{el conjunto de las partes abiertas~$\Gamma$ de~$X\times_S Y$}
tales que el morfismo~$\Gamma\to X$ inducido por~$\pr_1$ sea
\emph{radicial} y \emph{sobreyectivo}.
\end{corollaire}

\begin{proposition}
\label{IX.1.7}
Sea~$S_0$ el subpreesquema de~$S$ definido por un ideal cuasicoherente
formado por elementos nilpotentes, es decir, tal que~$S_0$ tenga el mismo
conjunto subyacente que~$S$. Entonces el funtor~$X\mto X\times_S S_0$ de
la categoría de los preesquemas étales sobre~$S$ a la categoría de los
preesquemas étales sobre~$S_0$ es una equivalencia de categorías.
\end{proposition}

El hecho de que este funtor sea plenamente fiel es una consecuencia
inmediata de~\Ref{IX.1.6}. El hecho de que sea esencialmente sobreyectivo
está contenido en~\Ref{I.8.3}. Se observará que, en la equivalencia
anterior, $X$ es de tipo finito, es decir, cuasifinito sobre~$S$
(respectivamente finito, es decir, un recubrimiento étale de~$S$), si y
solo si~$X_0$ satisface la condición análoga sobre~$S_0$; la misma
observación vale para la condición de separación. Estos hechos son
inmediatos y están contenidos también en~\Ref{IX.2.4} más abajo.

\begin{corollaire}
\label{IX.1.8}
Sea~$A$ un anillo local noetheriano completo de cuerpo residual~$k$.
Entonces el funtor~$B\mto B\otimes_A k$ es una equivalencia entre la
categoría de las álgebras finitas y étales sobre~$A$ y la categoría de las
álgebras finitas y étales sobre~$k$ (es decir, productos finitos de
extensiones finitas separables de~$k$).
\end{corollaire}

\begin{proposition}
\label{IX.1.9}
Para que~$X$ sea un \emph{recubrimiento étale} de~$S$, es decir, finito y
étale sobre~$S$, es necesario y suficiente que~$X$ sea $S$-isomorfo al
espectro de un álgebra~$\cal{A}$ sobre~$S$ que sea un módulo localmente
libre de tipo finito y tal que, para todo
\marginpar{231}
$s\in S$, $\cal{A}_s\otimes_{\cal{O}_s}\kres(s)$ sea un álgebra separable
sobre~$\kres(s)$ y, por tanto, en este caso un producto directo de
extensiones finitas separables de~$\kres(s)$.
\end{proposition}

Por último, el resultado siguiente es de naturaleza menos elemental, pues
es la conjunción de~\Ref{I.8.4} y del \emph{teorema de existencia para
haces en geometría algebraica} (EGA~III~5; \cf también \cite{IX.1},
teor\ptbl 3).

\begin{theoreme}
\label{IX.1.10}
Sean~$S$ el espectro de un anillo local noetheriano completo, $X$ un
$S$-esquema propio y~$X_0$ la fibra de~$X$ en el punto cerrado de~$S$ (de
modo que~$X_0$ es un subesquema cerrado de~$X$). Entonces el funtor de
restricción~$X'\mto X'\times_X X_0$ es una equivalencia entre la categoría
de los recubrimientos étales de~$X$ y la categoría de los recubrimientos
étales de~$X_0$.
\end{theoreme}
